\chapter{Calling convention}\label{app:abi}

The examples, the firmware and StrataOS use this convention. It is a software agreement, not an
architectural requirement.

\begin{table}[htbp]
\centering\small
\caption{Integer register usage}\label{tab:cc}
\begin{tabularx}{\linewidth}{@{}l l X l@{}}
\toprule
\textbf{Registers} & \textbf{ABI names} & \textbf{Use} & \textbf{Saved by}\\\midrule
\reg{x0} & \reg{zero} & constant zero & ---\\
\reg{x1} & \reg{ra} & return address (\insn{call}, \insn{jal ra}) & caller\\
\reg{x2} & \reg{sp} & stack pointer; the stack lives in the D-LS and grows down & callee\\
\reg{x3} & \reg{gp} & global pointer; the firmware requires \reg{gp} = \texttt{GPBASE} & ---\\
\reg{x4} & \reg{tp} & reserved & ---\\
\reg{x5}--\reg{x7}, \reg{x28}--\reg{x31} & \reg{t0}--\reg{t6} & temporaries & caller\\
\reg{x8}, \reg{x9}, \reg{x18}--\reg{x27} & \reg{s0}--\reg{s11} & saved registers (\reg{s0} is also \reg{fp}) & callee\\
\reg{x10}--\reg{x17} & \reg{a0}--\reg{a7} & arguments; \reg{a0}, \reg{a1} results & caller\\
\bottomrule
\end{tabularx}
\end{table}

\begin{itemize}
\item Bare-metal programs set \reg{sp} to 65536 (the top of the D-LS) at entry. The stack is kept 8-byte
  aligned. Under StrataOS the stack occupies D-LS \hex{E000}--\hex{FFFF} and a program's data D-LS
  \hex{0000}--\hex{9FFF}.
\item \reg{a0} at \insn{halt} is the program's exit status.
\item Capability registers have no general convention. The firmware conventions
  (Table~\ref{tab:sficonv}) are: \reg{c0} memory root, \reg{c2} console output, \reg{c3} device root,
  \reg{c9} block I/O buffer, \reg{c12}--\reg{c15} scratch. StrataOS programs receive their arena in
  \reg{c0}, endpoints to the servers in \reg{c1}--\reg{c4}, their message buffer in \reg{c5} and the
  clock in \reg{c7} (Table~\ref{tab:os-start}); \reg{c0}--\reg{c7} are preserved across context switches
  and \reg{c8}--\reg{c15} are not. A routine must not leave a live linear
  capability in a register its caller expects to be free, since a later write to it would trap.
\item Transfer tags: \reg{tg15} is used for console output by the examples; \reg{tg12}--\reg{tg15} are
  used by the firmware; \reg{tg10} by the StrataOS program library; \reg{tg14} by the StrataOS
  dispatcher, which a program must leave alone. Because tags affect only timing, sharing a tag is never incorrect.
\end{itemize}

\chapter{Trap causes}\label{app:causes}

\begin{table}[htbp]
\centering\small
\caption{Trap causes and trap values}\label{tab:causes}
\begin{tabularx}{\linewidth}{@{}r >{\sffamily}l X X@{}}
\toprule
\textbf{Code} & \textrm{\bfseries Name} & \textbf{Raised when} & \textbf{\csr{tval}}\\\midrule
1 & illegal & undefined opcode or function; a checked field with a value no instruction defines (Table~\ref{tab:fieldrules}); SYS function $>$ 2; CAP2 function $>$ 5; \insn{svc} in supervisor mode & the instruction word\\
2 & lsfault & load/store out of range or misaligned; transfer Local Store range out of range; bad tag-table entry; bad \insn{xmap} I-LS address & the Local Store address\\
3 & capinv & the named capability register or save slot is not live (empty, or revocable and its granule revoked); \insn{crevoke}/\insn{cundom} without a revoker; \insn{csend}/\insn{crecv} without a suitable endpoint & transfers and atomics: the far offset; capability and overlay instructions: the register or slot number\\
4 & perm & missing permission; transfer through an endpoint; \insn{tgetx} through device or core-control space; \insn{tget} through core-control space; \insn{catomic} outside T2/T3 & far offset (transfers, atomics) or register number\\
5 & bounds & far range outside the segment; overlay offset $\geq$ length & the far offset\\
6 & xalign & memory transfer operands or atomic offset not multiples of 8 & the far offset\\
7 & ifault & fetch from a misaligned \reg{pc} or one $\geq$ 65536 & \reg{pc}\\
8 & xmiss & \insn{xcall} target, or \insn{xret} far link, not mapped & the absolute far address; for \insn{xret}, the link (bit 63 set)\\
9 & pend & with CCR.PTRAP, a load, store, \insn{ltag} or \insn{linst} that would stall on an in-flight range; a \insn{tget} or \insn{tput} whose D-LS range overlaps a pending device read & the Local Store address\\
10 & capop & overwrite of a live linear capability; bad \insn{csplit} (beyond the end, an offset that is not a multiple of 8, or a revoker cut off a granule boundary) or \insn{cjoin}; \insn{cgive} outside the window (or in user mode without one), beyond node $N$ or to a halted core; \insn{csend} to such a node; bad \insn{cdom} source (not whole granules, revocable already, \ldots); an empty revoker; \insn{cmkep} beyond node $N$ & register or slot number, split offset, or node number\\
11 & priv & in user mode: \insn{csave}, \insn{crestore}, \insn{cdom}, \insn{crevoke}, \insn{cundom}, \insn{cmkep}; \insn{ctake}, \insn{cpoll}, \insn{tany} with the inbox bit; reading \csr{sscratch} or \csr{ipend}; writing \csr{status}, \csr{corebase}, \csr{corecount} or CSRs 19--24 & the instruction word\\
12 & halt & \insn{halt} in user mode (always to the supervisor) & 0\\
13 & stop & a stop request from core control, taken in user mode (always to the supervisor) & 0\\
14 & svc & \insn{svc} in user mode (always to the supervisor) & 0\\
\bottomrule
\end{tabularx}
\end{table}

Code 0 is never reported. Simulators print causes by these names. Causes 11--14 are always delivered
to the supervisor; the others go to the user handler when a user-mode program has installed one
(Section~\ref{sec:trapentry}).

\chapter{Control and status registers}\label{app:csrs}

\begin{table}[htbp]
\centering\small
\caption{Control and status registers}\label{tab:csrs}
\begin{tabularx}{\linewidth}{@{}r >{\ttfamily}l c c X@{}}
\toprule
\textbf{No.} & \textrm{\bfseries Name} & \textbf{S} & \textbf{U} & \textbf{Description}\\\midrule
0 & cycle & R & R & clock cycles since reset. Timing-visible; the emulator returns \csr{instret}\\
1 & instret & R & R & instructions retired (trapping instructions, \insn{svc} among them, do not retire; a supervisor \insn{halt} does)\\
2 & cause & RW & RW & cause of the last trap taken in the current mode (banked; Appendix~\ref{app:causes})\\
3 & epc & RW & RW & address of the trapping instruction, banked; \insn{eret} continues here\\
4 & tval & RW & RW & trap value, banked\\
5 & tvec & RW & RW & trap handler address, banked. Supervisor: 0 means a trap stops the core. User: 0 means every trap goes to the supervisor\\
6 & ccr & RW & RW & core control: bit 0 PTRAP (turn hazard stalls into \cause{pend} traps); the other bits are reserved: they hold what is written and have no effect\\
7 & coreid & R & R & this core's number, relative to \csr{corebase}\\
8 & ttag & RW & RW & software cache: D-LS address of the tag table\\
9 & tdata & RW & RW & software cache: D-LS address of the line storage\\
10 & tgeom & RW & RW & software cache geometry (Figure~\ref{fig:tgeom})\\
11 & ncores & R & R & \csr{corecount} if nonzero, otherwise the number of cores $N$\\
12 & scratch & RW & RW & free for software, shared by both modes\\
13 & stalls & R & R & cycles spent stalled. Timing-visible; the emulator returns 0\\
14 & takefrom & R & R & sender of the capability most recently taken by \insn{ctake} or \insn{crecv}, relative to \csr{corebase}\\
15 & corebase & RW & R & first physical core of the core window\\
16 & corecount & RW & R & size of the core window; 0 means no window\\
17 & reserved17 & --- & --- & reserved: reads as 0, writes are ignored\\
18 & status & RW & R & bit 0: 1 in supervisor mode (read-only); bit 1, SPP: \insn{eret} from supervisor mode enters supervisor mode if 1, user mode if 0\\
19 & ucause & RW & R & the user-mode copy of \csr{cause}\\
20 & uepc & RW & R & the user-mode copy of \csr{epc}\\
21 & utval & RW & R & the user-mode copy of \csr{tval}\\
22 & utvec & RW & R & the user-mode copy of \csr{tvec}: the program's trap handler, 0 for none\\
23 & sscratch & RW & --- & supervisor scratch register\\
24 & ipend & RW & --- & bit 0: a stop request is pending (set through core control; taken in user mode)\\
25 & xslot & RW & RW & bits 5:0: $\log_2$ of the overlay slot size; 0 selects exact mapping (Section~\ref{sec:xslot})\\
\bottomrule
\end{tabularx}
\end{table}

Columns S and U give the access in supervisor and user mode. R: read; RW: read and write; ---: the
access traps with \cause{priv}. In user mode, writing a register marked R traps with \cause{priv} if it
is \csr{status}, \csr{corebase}, \csr{corecount} or one of CSRs 19--22, and is ignored otherwise.
Reading an unassigned CSR number returns 0; writing a read-only or unassigned CSR has no effect. All
CSRs are zero at reset, except \csr{status}, which reads 3 (supervisor mode, SPP = 1).

\chapter{Encoding tables and instruction index}\label{app:index}

The primary opcode map is Table~\ref{tab:opmap} and the function assignments follow it in
Chapter~\ref{ch:encoding}. Table~\ref{tab:insindex} lists every instruction.

\input{gen-index}

\chapter{Glossary}\label{app:glossary}

\begin{description}[leftmargin=1em,style=nextline]
\item[As-if-synchronous] The rule that each transfer takes effect at its issue point; hardware hides the
  latency by stalling conflicting accesses.
\item[Atomic segment] Memory reachable only through an \perm{A} capability, created by \insn{catomic};
  accessed only with \insn{tamo}.
\item[Capability] An unforgeable (space, base, length, permissions, revocable flag, generation) descriptor that
  authorises far accesses.
\item[Code unit] A $2^{\csr{xslot}}$-byte block of code in far memory that slot mapping pages into an
  I-LS slot as a whole.
\item[Cluster store (T2)] Shared on-chip memory, 256 KiB in the prototype.
\item[Copyable] A capability with \perm{D}; moving it leaves the source in place.
\item[D-LS, I-LS] The data and instruction Local Stores of a core (T1d, T1i).
\item[Dead] Of a capability: valid and revocable, but its base's granule has been revoked since it
  was made; it behaves like an empty register.
\item[Drain] Wait for all of the core's in-flight transfers to complete.
\item[Endpoint] A capability in space 6 that names a node for \insn{csend} and \insn{crecv}.
\item[Dispatcher] The StrataOS kernel code on a worker core: it schedules the core's threads, switches
  contexts and demultiplexes the inbox.
\item[Doorbell] A stop request sent to a core so that its supervisor looks at its inbox.
\item[Event] Something a core waits for: a transfer completing or a gift arriving (\insn{tany}).
\item[Event port] A device port whose reads stay pending until the device has something to deliver:
  console input, console break, the timer's WAIT, the framebuffer controller's VSYNC.
\item[Far] Beyond the Local Stores: T2, T3, T4.
\item[Far link] A return link with bit 63 set: the far address of the instruction after an \insn{xcall}
  made from paged code, resolved by \insn{xret}.
\item[Framebuffer] The top 16~MiB of system memory, whose first rows the display scans out.
\item[Gift] A capability in an inbox, with its sender.
\item[Inbox] Where a node receives capabilities: for a core, one two-entry FIFO queue per sender; for the block device, one four-entry FIFO.
\item[Linear] A capability without \perm{D}; it can be moved but never copied or silently overwritten.
\item[Granule] 256~KiB of T3 or T2, the unit of revocation; the machine keeps a generation for each.
\item[Live] Of a capability: valid and, if it is revocable, of its base's granule's current
  generation.
\item[Node] A destination of \insn{cgive}: a core, or the block device (node $N$).
\item[Message] In StrataOS, a 4~KiB buffer capability lent to a server and given back with the result.
\item[Overlay] Code kept in far memory and loaded into the I-LS on demand; see \insn{xcall}.
\item[Pending device read] A read of an event port that has issued and not yet completed; its tag is
  busy.
\item[Platform table] The self-description of the machine (\texttt{STPD}) in device space.
\item[Publication rule] A hand-over (\insn{cgive}, \insn{csend}, \insn{cshare}, \insn{catomic},
  \insn{tamo}) waits until the core has no transfer to T2 or T3 in flight.
\item[Revocable] Of a capability: made from a range \insn{cdom} made revocable; it dies when the
  granule its base lies in is revoked.
\item[Revoker] The \perm{V} capability that \insn{cdom} returns; \insn{crevoke} and \insn{cundom} revoke
  its range. It can be cut at granule boundaries.
\item[Unit] A range of granules made revocable or revoked as one; a revocation kills the units it
  touches whole, and \insn{cjoin} never joins across units.
\item[Segment] The range of a far space that a capability names.
\item[Share] A non-writable copyable capability made by \insn{cshare}.
\item[Stop request] A doubleword written through core control; it makes a user-mode core trap to its
  supervisor with cause \cause{stop}.
\item[Supervised execution] Running a program in user mode in a window of cores set by
  \csr{corebase}/\csr{corecount}; its \insn{halt} traps to the supervisor.
\item[Supervisor call] The trap taken by \insn{svc} in user mode, cause \cause{svc}.
\item[Supervisor mode, user mode] The two privilege modes (Section~\ref{sec:modes}).
\item[System core] Core 0 under StrataOS: it runs only the kernel.
\item[System memory (T3)] Shared main memory, 64 MiB in the prototype; its top 16 MiB is the framebuffer
  window.
\item[Tag] One of 16 transfer identifiers used by \insn{twait}, \insn{tpoll} and \insn{tany}.
\item[Tick] The timer's unit: one microsecond.
\item[Timing-visible] State or behaviour whose value can depend on latency or scheduling.
\item[Transfer Unit (TU)] The per-core engine that performs \insn{tget}, \insn{tput} and \insn{tgetx}.
\end{description}

\chapter{Known limitations and future work}\label{app:limits}

\begin{description}[leftmargin=1em,style=nextline]
\item[No timer interrupt] The timer is a device, and preemption is a stop request that some core
  sends. An operating system that preempts dedicates a core, or part of one, to waiting for the timer;
  StrataOS dedicates core 0.
\item[One deadline] The timer has a single DEADLINE register. Software that needs several timeouts
  keeps a queue and gives the timer to one owner.
\item[Pending device reads] A core has at most four in flight. \insn{halt} waits for them, so a core
  must retire a read that may never complete before it halts. \insn{ltag} and \insn{linst} must not be
  used on a tag table or line that a pending device read is filling.
\item[Display] The geometry is fixed for a run and the only pixel format is XRGB8888. The FPGA builds
  have no video output.
\item[Supervisor code is never preempted] A stop request waits until the core runs in user mode. A
  supervisor that loops forever can be stopped only by resetting the machine.
\item[Revocation limits] Revocation works on whole 256~KiB granules. A granule's generation is 16
  bits wide, so a dead capability becomes live again once the generation of its base's granule has
  advanced through all 65,536 values while it is held. A revocation kills every unit it touches whole,
  also the parts outside the revoker's range.
\item[Block device and revocation] The block device performs a command through the capability it was
  given, even if that capability is revoked after the gift and before the command runs; the
  capability is returned dead.
\item[Serial Transfer Unit and bus] Each TU executes one transfer at a time, and the current far bus
  carries one transaction system-wide. This is sufficient for invariant I2 but limits bandwidth. A
  parallel TU must preserve the order of overlapping transfers and of device accesses.
\item[Inbox capacity] Each sender's queue at a core holds two gifts, so a sender that runs ahead of
  its receiver by more than two gifts waits. This is flow control, not a determinism hazard
  (Section~\ref{sec:determinism}).
\item[Atomic reads] An atomic segment can only be read with an atomic (e.g.\ \insn{tamo.add} of 0).
\item[Block device] Commands submitted beyond the queue depth are discarded silently; an empty
  completion queue reads as zeros; the return core in a command is a physical core number and is not
  translated by \csr{corebase}; a user-mode program reaches the device only through an endpoint for its
  node, and its supervisor must hand it the buffer that comes back.
\item[Cores] At most 15 cores: fifteen cores and the block device are the sixteen inbox senders that
  the masks of \insn{cpoll} and \insn{tany} describe.
\item[StrataOS] The limits of Table~\ref{tab:os-limits} apply: 11 threads per core, 1~MiB arenas,
  eight handles. A thread's \reg{c8}--\reg{c15} do not survive a context switch. An attachment is paired
  with the next message from the same core, so a process that sends attachments must be alone on its
  core or tolerate the pairing. A file being written is limited to 128~KiB, and the directory to 128
  contiguous files. The kernel is not preempted.
\item[Integer arithmetic] There are no 32-bit arithmetic forms, no signed high multiply, and no
  floating point.
\end{description}
